% !PRINTED 268
\weberpage{268}
%<<CONT
\[
\alpha^2 - \gamma = \delta^2,
\]
then $\alpha^2 > \delta^2$, and hence, if $\delta$ is taken positive,
\[
\alpha > \delta
\]
and
\[
27 D = 4 \delta^6 - (2\alpha\delta^2 - \beta)^2.
\]
By this, $D$ cannot be positive when $\beta$ is negative; for if $\beta$ is negative, then
\[
2\alpha\delta^2 - \beta > 2\alpha\delta^2 > 2\delta^3,
\]
and so $4\delta^6 - (2\alpha\delta^2 - \beta)^2$ is negative.

It is hereby directly proved that the criteria (5) and (6) say exactly the same thing.

According to the geometric interpretation given in \S.\,78, $\beta = 0$ means a surface of the third order which passes through the parabola $b = 0$, $\gamma = 0$, and which, so far as negative values of $a$ come into consideration, runs entirely within the region of space in which $D$ is negative.

In order to have before our eyes an example for the equation of the fifth degree, we consider the equation
\[
x^5 + x^2 + a = 0.
\]
We therefore have to set, in our general expressions of \S.\,73 or also of \S.\,74, (2), (3),
\[
a_0, a_1, a_2, a_3, a_4, a_5 =
\]
\[
1, 0, 0, 1, 0, a,
\]
and we obtain for the determinant of the Bezoutiant
\[
\begin{vmatrix}
-2a, & 0, & 0, & -5a \\
0, & \tfrac{6}{5}, & -5a, & 0 \\
0, & -5a, & 0, & -3 \\
-5a, & 0, & -3, & 0
\end{vmatrix}
\]
and for the discriminant
\[
D = 108 a + 3125 a^4.
\]
A chain of principal subdeterminants is
\weq{8}{D, \; 50a^3, \; \frac{-12a}{5}, \; -2a, \; 1.}
The discriminant is negative when
\weq{9}{0 < -a < \sqrt[3]{\frac{108}{3125}},}
otherwise positive, and so in particular positive for all positive $a$. One sees that when the discriminant is positive there are always two
%CONT>>
